def si d in k y f in k [X] ent la multiplicidad de a como raiz de f es el  numero {m in No : (X-alfa)^n|f}
c in {0,1,------,gr(f)} = num
siempre es un conjunto distinto de vpacio

obs num es un elemento de N0
tq(x-alfa)^num |f
(x-alfa)^num -1 no| f 
(x-alfa)^num + k no | f  k> 0 
obs:
num = 0 <-> alfa no es raiz 
obs:
(x-alfa)^num| f y num> 0 
x-alfa | (x-alfa)^num | f = b * p1^m1------ pr^mr 
x-alfa | p1
num< m1 
sol x -alfa = p1 

sea alfa1 .... alfa s a las raices de f en k 
(x-alfa1)^num1 , (x-alfa2)^num2 ...... , (x-alfa s)^num s|f
(x-alfa1)^num1 , (x-alfa2)^num2 ...... , (x-alfa s)^num s * g 

obs : alfa es raiz de f de mult d y de g de mult m , entonces que puedo decir de f*g 
f = (x-alfa)^num fbarrita fbarrita (alfa) distinc = 0 
g = (x-alfa)^m gbarrita gbarrita (alfa) distinc = 0 
fg = sumo exponentes fbarrita * gbarrita
obs = 0 distinto f+g = (x-alfa)^num fbarrita + (x-alfa)^m gbarrita
=(x-alfa)^m ((x-alfa)^num-m fbarrita + gbarrita)
h(alfa) = ((x-alfa)^num-m fbarrita(alfa) + g barrita(alfa)) distinct 0 
si f +g distinc 0 y mult(f,alfa) distinc = mult(g,alfa) , 
mult(f+g,alfa)

min{mult(f,alfa) mult(g,alfa)}

f(alfa) - q(alfa) * g(alfa) +r 

y alfa es raiz de g y f entonces alfa queda el resto ???revisar

obs: alfa raiz de f de mult num y de g de mult u
h = mcd(f,g)



sup f = alfa * p1^m1 * g con p irreducible m >= 1 
y p no divisible por g 

derivada:
d(fg) = d(f)*g + d(g)*f 

f' = (alfa^m*g) = alfa m p^m-1 p'* + alfa p^m * g' 
alfa p^m-1 ( mp'g + pg')
¿p| mp'g? NO 

lema = si alfa es una raiz de f de multiplicidad  m , entonces alfa es una raiz de f' de multiplicidad m-1 


c[x] talue x^7  -1 e^2pii/7
f' = 7x^6 

aca sacas que si alfa es raiz de f' es necesariamente raiz de f y no se cumple al reves
obs f in k{x} , k c C 
f divide factores repetidos -> mcd(f,f')=distict 1 
<-> f y f' no son coprimos 
f = p^m g n>=2 
p no | g f'= p^n-1(){esto igual a p no divide a r}

c no tiene ningun cuerpo que sea finito 


teorema fundamental del algebra todo polinomio no constante de c[X] tiene una raiz en c 
f in c[x] f= alfa * p1^m1 --------------- pr^mr 
un pol en complejo se factoriza en solo sus raices nada que sea del factor g =1 


x^2 +1 in R[x] no tiene ninguna raiz in R 
x^2 -2 om Q[x]                       in Q
x^2 + x +1 in Z2[x]

obs : si yo tengo un polinomio de grado impar si o si corta al eje x en algun punto osea tiene raiz, entonces hay alfa in R talque f(alfa) = 0 , asi que x-alfa|f 
rta  F = x- alfa
si f es irreducible en R[x] de frado par, entonces f tiene grado 2 y no tiene raices en R 


f = x^2 + ax +b con a^2 -4b < 0 osea discriminante negativo osea no real 

-a +- raiz cuadrada de (a^2 -4b)/ 2


en reales las raices son de grado 1 la x o son de grado 2 con discriminante <0 

obs: sobore q hay polinomios irreducibles de todos los grados 

sea f in R[x] de grado par irreducible sea alfa1, --- alfar ñas raices complejas  de f . ninguno de ellos esta en R s